\subsection{Four positive helicities and two negative}\label{sec:eight-two-minus}

Consider $\M(1_e,2_E,3^+,4^+,5^+,6^+,7^-,8^-)$. Split the charged
line at its middle electron propagator. Write $D(ijk|lmn)$ for the
contribution with the three photons $i,j,k$ at the electron end and
$l,m,n$ at the positron end, summed over their orders on each side.
There are $\binom{6}{3}=20$ such partitions, each accounting for
$3!\,3!=36$ of the 720 diagrams. The same function $D$ gives both
pole forms.

\paragraph{Physical poles.}
Use $p_1$ for every polarization reference. In this gauge every orbital
factor $x_i(p_1)$ in the electron's pair current vanishes. Let
$(a,b,c,d)$ run over all 24 permutations of $(3,4,5,6)$ and $(h,g)$
over both orders of $(7,8)$. The number of negative photons at the
electron end distinguishes the three terms in the result.
\begin{equation}\label{eq:eight-two-physical}
 \M^{\phys}_{++++--}=-(\sqrt{2}e)^6\sum_{S_4\times S_2}
 \left[\frac{D(abc|dhg)+D(ahg|bcd)}{12}
       +\frac{D(abh|cdg)}{4}\right].
\end{equation}
The three choices account for four, four and twelve partitions,
respectively. The divisors remove repetitions within each unordered triple.

\paragraph{Spurious poles.}
Use photon $h$ as the angle reference of every positive photon. The
two negatives serve as each other's square references. All six
references are then external photons. Reversing the line interchanges
the triples with a minus sign in our all-incoming convention. With
$P_{12}$ exchanging the massive momenta and spin labels, the result is
\begin{equation}\label{eq:eight-two-spurious}
 \M^{\spur}_{++++--}=-(\sqrt{2}e)^6(1-P_{12})\sum_{S_4\times S_2}
 \left[\frac{D(abc|dhg)}{12}+\frac{D(abh|cdg)}{8}\right].
\end{equation}
At fixed $h$ the first term generates the four partitions with three
positives on the left. The second generates the six with $h$ on the
left and $g$ on the right. Reversal supplies the other ten. The sum
over $(h,g)$ averages the two gauges of the complete amplitude. This
accounts for the additional factor of two in the divisors relative to
summing only the positives at fixed $h$. The only spurious poles come
from the external polarization brackets.

\paragraph{The middle propagator.}
For a photon $i$ write $\alpha_i$ and $\beta_i$ for its angle and
square polarization spinors, and $n_i$ for the polarization denominator.
The following choices specify every bracket in $D$.
\[
\begin{array}{c|ccc|ccc}
 &\multicolumn{3}{c|}{\text{physical}}&\multicolumn{3}{c}{\text{spurious}}\\
 &\alpha_i&\beta_i&n_i&\alpha_i&\beta_i&n_i\\ \hline
 i^+&p_1|i]&i&\tau_{1i}&r_i&i&\langle r_i i\rangle\\
 i^-&i&p_1|i\rangle&\tau_{1i}&i&r_i&[r_i i]
\end{array}
\]
In the physical columns a projected spinor at the left of a bracket
inserts $p_1$ with a plus sign, and one at the right inserts it with
a minus sign. For example,
$\langle X\alpha_i\rangle=-\langle X|p_1|i]$ and
$\langle\alpha_i X\rangle=[i|p_1|X\rangle$ for $i^+$.
The rule for a projected square spinor follows by exchanging brackets.

Let $J_t(ij)$ be the open two-photon spinor at endpoint $t$ from
Eq.~\eqref{eq:seven-two-pair}, using the references of the chosen
pole form. Summing over the photon that touches the middle
propagator gives the two left spinors
\begin{align}
 \langle L|&=\sum_{\rm cyc(i,j,k)}
   \frac{[J_1(ij)\,\beta_k]}{n_k}\langle\alpha_k|,&
 [L|&=\sum_{\rm cyc(i,j,k)}
   \frac{\langle J_1(ij)\,\alpha_k\rangle}{n_k}[\beta_k|.\nonumber
\end{align}
The corresponding right spinors are
\begin{align}
 |R\rangle&=\sum_{\rm cyc(l,m,n)}
  |\alpha_l\rangle\frac{[\beta_l\,J_2(mn)]}{n_l},&
 |R]&=\sum_{\rm cyc(l,m,n)}
  |\beta_l]\frac{\langle\alpha_l\,J_2(mn)\rangle}{n_l}.\nonumber
\end{align}
With $Q=p_1+p_i+p_j+p_k$ and $\tau_{1ijk}=Q^2-m^2$, sewing the two currents gives
\begin{equation}\label{eq:eight-two-bridge}
 D(ijk|lmn)=
 \frac{m\big(\langle LR\rangle+[LR]\big)
       +\langle L|Q|R]+[L|Q|R\rangle}{\tau_{1ijk}} .
\end{equation}
The mass connects equal chiralities and the momentum connects opposite
chiralities. The middle propagator retains its full virtuality.
Equivalently one can multiply its numerator into either three-photon
current before summing the photon adjacent to it. The reducer compares
these expressions separately for each partition in the permutation sum.

\paragraph{Numerical checks.}
Both constructions satisfy all 62 physical residues and the fixed-angle
$E^{-4}$ bound. For the two paper formulas, the largest relative
discrepancies from the signed sum of 720 Feynman diagrams at 100,000
points were $8.45\times10^{-11}$ in both gauges. At a separate 100,000
points, the squared amplitudes with the massive spins summed agreed
with polarized MadGraph to $5.87\times10^{-12}$ and $1.12\times10^{-10}$
in the physical and spurious forms, respectively. Both comparisons met
their tolerances throughout in native arithmetic. The automatically
rearranged summands also passed both samples.
